Korean financial privacy detection must address jurisdiction-specific categories, noncanonical identifiers and contextual attributes. We introduce \textbf{Kiii-Kiii} (dataset \kiii{}), a regulation-grounded benchmark with 36 categories: 24 with statutory grounds and 12 additional quasi-identifier and profiling categories. The taxonomy separates detection policy from legal redaction obligations. Its synthetic preview contains 1,440 documents and 218,664 spans across ten document types, with code-generated identifiers and four levels of surface-form variation. A matched full-context versus local-window protocol fixes target regions and output budgets. Across 12 language models and three frozen baselines, strict exact-span extraction is limited by output-contract failures and missed spans. The best LLM strict F1 is 7.97. A fixed post-hoc itemwise diagnostic reaches 31.58 F1 while retaining all gold and penalizing invalid items; it measures observable recovery, not latent ability. Local context yields higher diagnostic F1 in all ten completed pairs; strict scores favor local in nine. Archived-response analysis distinguishes empty outputs, core-ownership violations, and category or boundary mismatches, showing why compliance and detection require separate reporting. The preview is released under CC BY-NC 4.0, with a formal release to follow.

